\section*{Chapter \ref{ch:rc:inflation-derivatives} --- Inflation Derivatives}

\begin{solution}{exo:rc:inflation-derivatives:1}
The forward index ratio is $1.0328^{10} = 1.3809$; the real discount factor is
$P_N(0,10)\times1.3809 = 0.6771\times1.3809 = 0.9349$.
\end{solution}

\begin{solution}{exo:rc:inflation-derivatives:2}
The first year's ratio $I(T_1)/I(0)$ has a known denominator: it is a zero-coupon payment,
whose expectation under the payment measure is the forward index ratio exactly.
\end{solution}

\begin{solution}{exo:rc:inflation-derivatives:3}
By 5\% (the cap binds); with $[0\%,2.5\%]$, by 2.5\%.
\end{solution}

\begin{solution}{exo:rc:inflation-derivatives:4}
The expectation of the year's ratio is the expected real bond price $P_R(T,S)$ at the
start of the year, over the known nominal one. Real-rate uncertainty makes the expected
bond price (under the measure used) lower than its forward, through the drift $\mu(T)$ that
stochastic real rates imply: the adjustment is negative and grows with $T$. A positive
$\rho$ lowers the real rate's drift under the nominal measure (the quanto term), which
raises the expected bond price and offsets part of it.
\end{solution}

\begin{solution}{exo:rc:inflation-derivatives:5}
The May index: from January to June the US pattern adds about one point to the forward
relative to the smooth path, while at December the pattern has summed to zero and the
forward is back on the smooth path. The swap fixing on May is therefore worth more.
\end{solution}

\begin{solution}{exo:rc:inflation-derivatives:6}
Each year's capped and floored change multiplies the amount built up so far: the payoff is
a product of options, whose value depends on the joint distribution of all years' changes
and on the path, not a sum of options on single years.
\end{solution}

\begin{solution}{exo:rc:inflation-derivatives:7}
$-0.0182$ per unit when $\sigma_I$ goes from 1.5\% to 2.5\%, $-0.0061$ when $\sigma_R$ goes from 0.80\%
to 0.90\%: the leg is short volatility (the cap at 5\% is worth more to the payer as volatility
rises than the floor at 0\%).
\end{solution}

\begin{solution}{exo:rc:inflation-derivatives:8}
The zero-coupon rate compounds one ratio of forward indices; the year-on-year rate averages
the expectations of annual ratios, each shifted by its convexity adjustment and weighted by
nominal discount factors. With stochastic real rates the two differ: 3.22\% against 3.28\% at
ten years here.
\end{solution}

\begin{solution}{pb:rc:inflation-derivatives:1}
\textbf{1.} 0.8356 by simulation; exactly $P_R(0,20) = 0.8345$ (the difference is sampling
error).
\textbf{2.} 0.7863.
\textbf{3.} 0.6371.
\textbf{4.} It is the index with a short cap at 5\% and a long floor at 0\% on each year's change,
compounded; with expected inflation near 3.3\%, the cap is closer and costs more than the
floor is worth.
\textbf{5.} An uncapped hedge overhedges: it pays inflation above 5\% the liabilities never
incur, so the scheme holds an unhedged long-inflation, long-volatility position worth about
$0.8356-0.7863 = 0.049$ per unit.
\textbf{6.} $-0.0182$ per unit.
\textbf{7.} $-0.0061$ per unit.
\textbf{8.} Holding LPI liabilities and receiving LPI from the swap leaves it hedged; the
counterparty paying LPI is long volatility (short the leg), since the leg loses value as
volatility rises.
\textbf{9.} So that the difference measures the parameter's effect, not the noise of two
independent simulations.
\textbf{10.} Year-on-year caps and floors on the index, and index options where they trade;
correlation and path dependence remain.
\textbf{11.} 5\%.
\textbf{12.} 14.2\% on the year.
\textbf{13.} The hedge's excess over the liabilities was a large gain unrelated to the scheme's
obligations: a position, not a hedge, that would be lost if inflation reversed.
\textbf{14.} From 2030 the RPI is computed like CPIH, which has run below it; forward RPI
beyond 2030 falls, and so do the values of RPI-linked LPI legs.
\textbf{15.} Without a floor the leg is the index with a cap only: cheaper still than the
$[0\%,5\%]$ leg, and the deflation risk sits with the scheme's members.
\textbf{16.} Institutions with the opposite exposure: issuers of LPI-linked assets (some
infrastructure, property leases), insurers writing annuity business, and dealers who warehouse
the volatility.
\textbf{17.} It drives the convexity and the joint distribution of annual changes, and no liquid
instrument quotes it; its choice moves the leg's value.
\textbf{18.} The value of the capped liabilities, the mismatch with any uncapped hedges, the
volatility and correlation sensitivities, and a 2022-like scenario.
\textbf{19.} Named result: \emph{the LPI swap}: the uncapped leg is worth 0.8345 per unit (0.8356
simulated), the $[0\%,5\%]$ LPI leg 0.7863, and one point more index volatility lowers the LPI leg by
0.0182.
\textbf{20.} It turns the promise into an option-laden payoff that is worth less than the index
and loses value as inflation becomes more volatile.
\end{solution}

\begin{solution}{iq:rc:inflation-derivatives:1}
Nominal money is the domestic currency, a unit of the basket is a foreign currency, and
the price index is the exchange rate. Real bonds are foreign bonds; the forward index is
a forward exchange rate $I(t)P_R/P_N$; real rates seen from the nominal measure acquire a
quanto drift. Cross-currency models become inflation models.
\iqlookfor{the mapping and its consequences.}
\end{solution}

\begin{solution}{iq:rc:inflation-derivatives:2}
Each payment is a ratio of two future index values paid at the later date; its
expectation under the payment measure is the forward ratio times a convexity adjustment
that depends on real-rate volatility and correlations, which the zero-coupon curve does
not contain.
\iqlookfor{ratio of future values and the convexity adjustment.}
\end{solution}

\begin{solution}{iq:rc:inflation-derivatives:3}
Swaps and bonds reference a given month's index, not seasonally adjusted; the forward
between annual points must carry the monthly pattern. Estimate seasonal factors from the
index history (complete years), keep the annual swap points, and bend the path between them;
check stability of the factors over time.
\iqlookfor{monthly reference and historical estimation.}
\end{solution}

\begin{solution}{iq:rc:inflation-derivatives:4}
Annual increases equal to the index's change floored and capped (in the UK, $[0\%,5\%]$ or
$[0\%,2.5\%]$), compounded. Hedges are path-dependent options on inflation, illiquid, with
volatility and correlation exposure; uncapped hedges overhedge in high-inflation years.
\iqlookfor{path dependence and the 2022 lesson.}
\end{solution}

\begin{solution}{iq:rc:inflation-derivatives:5}
Observable: nominal and real curves (swaps, linkers), and, less liquidly, index
volatilities from year-on-year and zero-coupon options. Weakly observable or not: the real
rate's mean reversion and volatility, and the correlations between nominal rates, real
rates and the index; they are set from history or implied from few options.
\iqlookfor{which parameters are marked, not calibrated.}
\end{solution}

\begin{solution}{iq:rc:inflation-derivatives:6}
The drift of the index ($n-r^R-\frac12\sigma_I^2$), the quanto drift of the real rate
($-\rho\sigma_R\sigma_I$), the real curve's construction from the swaps (compounding,
lag), time-grid alignment with the payment dates, and discounting with the nominal curve;
test that the simulated index times the nominal discount factor reprices $P_R(0,T)$ at each
maturity.
\iqlookfor{drifts, measures and a direct repricing test.}
\end{solution}
